% ===========================================================================
% Your resume. Replace every placeholder below, then delete this banner.
%
%   Requires XeLaTeX.  Overleaf: Menu -> Settings -> Compiler -> XeLaTeX
%   Build:  latexmk -xelatex -outdir=build main.tex
%   Check:  make check          (fails while placeholders remain)
%
%   The vocabulary is seven primitives -- read docs/GRAMMAR.md before writing.
%   examples/academic.tex and examples/industry.tex are two full compositions;
%   neither is "the" shape of a resume. Invent the sections you actually need.
%
%   Class options:
%     fontset    = commercial (included config) | your own set -- docs/FONTS.md
%     pagefooter = false | true  ("Name . 2/3" centred at the foot)
%     icons      = true | false  (contact icons)
% ===========================================================================
\documentclass[fontset=commercial,pagefooter=false]{yuan2resume}

\cvname{Your Name}
\cvposition{Your Job Title}

% The contact block. Items print in declaration order, so to reorder them you
% reorder these lines; to move an item to the other row, move its line into
% the other \cvcontactrow. Every item must be inside an explicit row.
% Icons are semantic slots (\cvIconPhone, \cvIconScholar, ...), so you never
% name a glyph here -- and icons=false in the class options drops them all.
\cvcontactrow{
  \cvcontact[\cvIconPhone]{+1 (234) 567-8900}
  \cvemail{your.email@example.com}
}
\cvcontactrow{
  \cvlink[\cvIconScholar]{https://scholar.google.com/citations?user=XXXX}{Scholar}
  \cvlink[\cvIconGitHub]{https://github.com/you}{github.com/you}
  \cvlink[\cvIconSite]{https://example.com}{example.com}
}
% Other slots: \cvIconLinkedIn, \cvIconLocation, or pass any fontawesome5
% command directly -- \cvlink[\faTwitter]{...}{...}

\begin{document}
\maketitle

\begin{cvtwocolumn}

\begin{cvsection}{Experience}
  \cvpair{\textbf{Employer Name}}{City, Country}
  \cvpair{\textbf{Your Job Title}, Team}{2022 -- \textit{Present}}
  \begin{cvitemize}
    \item What you owned, and the number that shows it mattered.
  \end{cvitemize}
  % cventry nests a project under the role. It indents from the left only, so
  % the date below still lands on the same right edge as "City, Country".
  \begin{cventry}
    \cvpair{\textit{Project Name}}{2023 -- 2024}
    \begin{cvitemize}
      \item What changed, measured.
    \end{cvitemize}
  \end{cventry}

  \cvgap  % <- separates sibling entries. Never write \vspace here.

  \cvpair{\textbf{Earlier Employer}}{City, Country}
  \cvpair{\textbf{Earlier Title}}{2019 -- 2022}
  \begin{cvitemize}
    \item One line per thing worth an interview question.
  \end{cvitemize}
\end{cvsection}

\begin{cvsection}{Education}
  \cvpair{\textbf{University Name}}{City, Country}
  \cvpair{\textit{Your Degree}}{2015 -- 2019}
\end{cvsection}

\begin{cvsection}{Skills}
  \textbf{Languages}: {\CodeFont{Python}}, {\CodeFont{Go}} \\
  \textbf{Tools}: name only what you would defend in an interview
\end{cvsection}

\end{cvtwocolumn}
\end{document}
